\section{Limitations}

Despite its advantages, our method exhibits several limitations that require further research.

\textbf{Implicit Temporal Representation}. Although timestamps can be added to triplet`s attributes, their reliance on explicit conversion to plain text (for LLM prompting)  creates inefficiencies. Due to the "Lost in the Middle" problem~\cite{liu-etal-2024-lost} it leads to potential loss of critical contextual data, thereby compromising overall search accuracy.

\textbf{Simplified Ontology Structure}. The current memory design offers limited characteristics for indexing and filtering information, resulting in suboptimal query performance and reduced effectiveness of Question-Answering (QA) algorithms.

\textbf{Ambiguous Entity Definitions}. Object vertices lack formal entity definitions, causing difficulties in resolving ambiguities during QA pipeline execution. Consequently, searches involving polysemous terms require extensive traversals through the memory graph, leading either to incomplete responses or false positives.

\textbf{Lack of Semantic Deduplication}. Current duplicate detection mechanism using only exact string comparisons rather than semantic equivalence. As such, synonymous triplets may be unnecessarily replicated, increasing storage demands, slowing down retrievals and complicating updates, particularly when pruning obsolete vertices and edges.
   
Addressing these issues represents key areas for future research aimed at improving both scalability and robustness of personalized Knowledge Graph-based QA systems.


% TODO

% процедурная фильтрация
% неявная темпоральна компонента
% слабая спецификация плана поиска
% не учитывается формат входящего текста при парсинге
% наивный способ отбора сопоствленных вершин при генерации clue-вопросов
% наивная проверка добавляемой информации в граф на дублирование